\section{Acciones y recomendaciones para la puesta en práctica}

\subsection{Lista de acciones a corto plazo}

\begin{table}[H]
\centering
\begin{tabular}{p{3cm}p{6cm}}
\toprule
Ventana de tiempo & Entregables \\
\midrule
Week 1 & Completar data review + skill segmentation; establecer prompt log template \\
Week 2 & Launch evaluation dashboard, registrar multi-metric drift; definir incident grade-book \\
Week 3 & Activar un emergent drill: LLM plan + skill exec y governance review \\
\bottomrule
\end{tabular}
\caption{Acciones recomendadas a corto plazo.}
\end{table}

\begin{importantbox}{Puntos clave para poner en práctica las acciones}
\begin{itemize}
    \item Prompt log + LLN plan deben llevar asociados version/temperature/expected skill;
    \item cada emergent drill debe contar con review notes y dar lugar a lessons learned;
    \item Incident grade-book registra el trigger, la redirect action y el guardrail change posterior.
\end{itemize}
\end{importantbox}

\subsection{Agenda de investigación a largo plazo}

Se recomienda organizar las líneas a largo plazo en forma de quarterly thesis: 1) scale-law + hierarchical optimizer; 2) LLM-driven skill discovery; 3) attention interpretability for subgoal.

\begin{knowledgebox}{Ejemplos de investigación a largo plazo}
\begin{itemize}
    \item Project Alpha: explorar patrones de hierarchical attention;
    \item Project Beta: mapear las curvas de scale law a un optimizer schedule;
    \item Project Gamma: construir un emergent dashboard que apoye las decisiones del governance panel.
\end{itemize}
\end{knowledgebox}

\subsection{Checklist automation}

Se automatiza la action checklist de corto plazo; por ejemplo, con Github Actions se verifica cada semana que el prompt log esté sincronizado y con dashboards se muestra el multi-metric drift. Así se reduce la verificación manual repetitiva y los ingenieros pueden concentrar su esfuerzo en la incident response.

\begin{importantbox}{Recordatorios de la checklist automatizada}
\begin{itemize}
    \item una actualización del prompt activa el workflow \texttt{check\_prompt\_versions};
    \item cuando un drift alert alcanza el threshold, se activa \texttt{incident-drill};
    \item \texttt{summary.md} genera cada semana un resumen de lessons para que lo revise el governance panel.
\end{itemize}
\end{importantbox}

\subsection{Resumen de la sección}

La sección de acciones recomendadas convierte la teoría en un timeline: los drills de corto plazo cierran el ciclo con rapidez, mientras que los proyectos de largo plazo convierten los emergent insights en research output verificable.
